% ============================================================
%  CHAPTER 4 — Results and Discussion
%  Typical length: 20–35 pages
%  Cover: present results, interpret findings,
%         compare against benchmarks or prior work.
% ============================================================

% Section are made as examples to show how to write — you have to make sections/subsections as per your need.


\chapter{Results and Discussions}
\label{chap:results}

[Section are made as examples — make your own sections/subsections headings as per your need.]

% TODO: Brief overview paragraph — what does this chapter cover?

% ---- 4.1 [First set of results] ----------------------------
\section{[First Result / Analysis]}
\label{sec:result1}

% TODO: Present your first key result.
% Lead with the figure or table, then explain what it shows.
% Interpret — don't just describe.

[Write your first result here.]

% ---- 4.2 [Second set of results] ---------------------------
\section{[Second Result / Analysis]}
\label{sec:result2}

[Write your second result here.]

\subsection{[Sub-result]}

[Write here.]

% ---- 4.3 [Comparative / Validation study] ------------------
\section{[Comparison / Validation]}
\label{sec:validation}

% TODO: Compare your results to prior work, analytical predictions,
% or alternative approaches. Discuss agreement and discrepancies.

[Write your validation or comparison here.]

% ---- 4.4 Discussion ----------------------------------------
\section{Discussion}
\label{sec:discussion}

% TODO: Synthesise all results. What do they mean collectively?
% What are the strengths and limitations of your findings?
% How do they relate to the objectives stated in Chapter 1?

[Write your overall discussion here.]


% Uncomment the below lines (if you have used AI for the contents of this chapter) to include an AI use disclosure at the end of the chapter. You can modify the wording as needed, but be transparent about how you used AI tools in writing this chapter.

% \vspace{1.0cm}
% \noindent\rule{\textwidth}{0.4pt}
% \vspace{0.2cm}

% \noindent\textit{Note on AI assistance:} An AI language model was used in this chapter as a collaborative tool for [describe how — e.g., interpreting results, discussion of findings, figure preparation]. All analytical decisions and interpretations are the author's own.
